\chapter*{练习}
\addcontentsline{toc}{section}{练习}

\section*{练习 5.1}
\addcontentsline{toc}{subsection}{练习 5.1}

计算下列函数的导数 $ f'(x) $：

\begin{equation*}
f(x) = \log(x^4) \sin(x^3)
\end{equation*}

\medskip\hrule\medskip

\section*{练习 5.2}
\addcontentsline{toc}{subsection}{练习 5.2}

计算 Logistic Sigmoid 函数：

\begin{equation*}
f(x) = \frac{1}{1 + \exp(-x)}
\end{equation*}

的导数 $ f'(x) $。

\medskip\hrule\medskip

\section*{练习 5.3}
\addcontentsline{toc}{subsection}{练习 5.3}

计算函数：

\begin{equation*}
f(x) = \exp\left(-\frac{1}{2\sigma^2} (x - \mu)^2\right)
\end{equation*}

的导数 $ f'(x) $，其中 $ \mu, \sigma \in \mathbb{R} $ 为常数。

\medskip\hrule\medskip

\section*{练习 5.4}
\addcontentsline{toc}{subsection}{练习 5.4}

计算函数 $ f(x) = \sin(x) + \cos(x) $ 在 $ x_0 = 0 $ 处的泰勒多项式 $ T_n $（$ n = 0, \ldots, 5 $）。

\medskip\hrule\medskip

\section*{练习 5.5}
\addcontentsline{toc}{subsection}{练习 5.5}

考虑下列函数：

\begin{equation*}
f_1(\boldsymbol{x}) = \sin(x_1) \cos(x_2), \quad \boldsymbol{x} \in \mathbb{R}^2
\end{equation*}

\begin{equation*}
f_2(\boldsymbol{x}, \boldsymbol{y}) = \boldsymbol{x}^\top \boldsymbol{y}, \quad \boldsymbol{x}, \boldsymbol{y} \in \mathbb{R}^n
\end{equation*}

\begin{equation*}
f_3(\boldsymbol{x}) = \boldsymbol{x} \boldsymbol{x}^\top, \quad \boldsymbol{x} \in \mathbb{R}^n
\end{equation*}

a. $ \frac{\partial f_i}{\partial \boldsymbol{x}} $ 的维度分别是什么？ b. 计算对应的雅可比矩阵（Jacobians）。

\medskip\hrule\medskip

\section*{练习 5.6}
\addcontentsline{toc}{subsection}{练习 5.6}

求 $ f $ 关于 $ \boldsymbol{t} $ 的导数以及 $ g $ 关于 $ \boldsymbol{X} $ 的导数，其中：

\begin{equation*}
f(\boldsymbol{t}) = \sin(\log(\boldsymbol{t}^\top \boldsymbol{t})), \quad \boldsymbol{t} \in \mathbb{R}^D
\end{equation*}

\begin{equation*}
\begin{adjustbox}{max width=\linewidth}
g(\boldsymbol{X}) = \operatorname{tr}(\boldsymbol{A} \boldsymbol{X} \boldsymbol{B}), \quad \boldsymbol{A} \in \mathbb{R}^{D \times E}, \; \boldsymbol{X} \in \mathbb{R}^{E \times F}, \; \boldsymbol{B} \in \mathbb{R}^{F \times D}
\end{adjustbox}
\end{equation*}

其中 $ \operatorname{tr} $ 表示矩阵的迹。

\medskip\hrule\medskip

\section*{练习 5.7}
\addcontentsline{toc}{subsection}{练习 5.7}

利用链式法则计算下列函数的导数 $ \frac{\mathrm{d}f}{\mathrm{d}\boldsymbol{x}} $。给出每个偏导数的维度，并详细描述你的推导步骤。

a.

\begin{equation*}
f(z) = \log(1 + z), \quad z = \boldsymbol{x}^\top \boldsymbol{x}, \quad \boldsymbol{x} \in \mathbb{R}^D
\end{equation*}

b.

\begin{equation*}
\begin{adjustbox}{max width=\linewidth}
f(\boldsymbol{z}) = \sin(\boldsymbol{z}), \quad \boldsymbol{z} = \boldsymbol{A} \boldsymbol{x} + \boldsymbol{b}, \quad \boldsymbol{A} \in \mathbb{R}^{E \times D}, \; \boldsymbol{x} \in \mathbb{R}^D, \; \boldsymbol{b} \in \mathbb{R}^E
\end{adjustbox}
\end{equation*}

其中 $ \sin(\cdot) $ 逐元素作用于 $ \boldsymbol{z} $。

\medskip\hrule\medskip

\section*{练习 5.8}
\addcontentsline{toc}{subsection}{练习 5.8}

计算下列函数的导数 $ \frac{\mathrm{d}f}{\mathrm{d}\boldsymbol{x}} $。详细描述你的推导步骤。

a. 使用链式法则。给出每个偏导数的维度。

\begin{equation*}
f(z) = \exp\left(-\frac{1}{2} z\right)
\end{equation*}

\begin{equation*}
z = g(\boldsymbol{y}) = \boldsymbol{y}^\top \boldsymbol{S}^{-1} \boldsymbol{y}
\end{equation*}

\begin{equation*}
\boldsymbol{y} = h(\boldsymbol{x}) = \boldsymbol{x} - \boldsymbol{\mu}
\end{equation*}

其中 $ \boldsymbol{x}, \boldsymbol{\mu} \in \mathbb{R}^D $，$ \boldsymbol{S} \in \mathbb{R}^{D \times D} $。

b.

\begin{equation*}
f(\boldsymbol{x}) = \operatorname{tr}(\boldsymbol{x} \boldsymbol{x}^\top + \sigma^2 \boldsymbol{I}), \quad \boldsymbol{x} \in \mathbb{R}^D
\end{equation*}

这里 $ \operatorname{tr}(\boldsymbol{A}) $ 是 $ \boldsymbol{A} $ 的迹，即对角线元素之和 $ \sum_i A_{ii} $。提示：显式写出外积。

c. 使用链式法则。给出每个偏导数的维度。不需要显式计算偏导数的乘积。

\begin{equation*}
\boldsymbol{f} = \tanh(\boldsymbol{z}) \in \mathbb{R}^M
\end{equation*}

\begin{equation*}
\boldsymbol{z} = \boldsymbol{A} \boldsymbol{x} + \boldsymbol{b}, \quad \boldsymbol{x} \in \mathbb{R}^N, \; \boldsymbol{A} \in \mathbb{R}^{M \times N}, \; \boldsymbol{b} \in \mathbb{R}^M
\end{equation*}

这里，$ \tanh $ 逐分量作用于 $ \boldsymbol{z} $。

\medskip\hrule\medskip

\section*{练习 5.9}
\addcontentsline{toc}{subsection}{练习 5.9}

对于可微函数 $ p, q, t $，我们定义：

\begin{equation*}
g(z, \nu) := \log p(x, z) - \log q(z, \nu)
\end{equation*}

\begin{equation*}
z := t(\epsilon, \nu)
\end{equation*}

利用链式法则，计算梯度：

\begin{equation*}
\frac{\mathrm{d}}{\mathrm{d}\nu} g(z, \nu)
\end{equation*}
